\section{The precise scope of a positive Rankin comparison}
\label{sec:s5-rankin-scope}

This section tests a proposed replacement of the regular
\(S_5\)-representation by the known positive Rankin functions.
It proves the finite local comparisons, including at ramified
places, and an exact limitation of the resulting leading
near-one score.  It does not compare every possible full
Euler--gamma variational expression, and does not reduce
the arithmetic gate in MA-RB-008.

\paragraph{Review status.}
The complete post-writing author reading and two independent
whole-proof readings have passed.  The source-hypothesis
checks below have their own recorded scope.

Retain the actual extension \(L/\mathbb Q\), its discriminant
\(\Delta\), and the representations and actual order matrix
of MA-RB-006.  Write \(\operatorname{Reg}=\mathbb C[S_5]\).
For a finite Artin representation \(T\), let \(N(T)\)
denote its finite Artin conductor.  At a prime \(p\), use
its actual inertia group \(I_p\) and a lift \(\phi_p\)
of arithmetic residue Frobenius, and set
\[
 t_{T,p}(j)=
       \operatorname{tr}\bigl(\phi_p^j\mid T^{I_p}\bigr),
 \qquad
 E_T(\sigma)=\sum_p\sum_{j\ge1}
             t_{T,p}(j)(\log p)p^{-j\sigma}
              \quad(\sigma>1).
 \tag{MA-RB.150}
\]
For \(T=\operatorname{Reg}\), this is
\(-\zeta_L'/\zeta_L(\sigma)\).

\paragraph{MA-RB-011 (local positivity and the five-factor score).}
Let \(U\) be any finite-dimensional complex representation
of \(S_5\), of dimension \(d\).  Then
\[
 \begin{gathered}
 \log N(\operatorname{End}U)
                       \le \frac{d^2}{120}\log\Delta,\\
 t_{\operatorname{End}U,p}(j)
                  \ge\frac{d^2}{120}\,t_{\operatorname{Reg},p}(j)
                           \qquad(p,\ j\ge1).
 \end{gathered}
 \tag{MA-RB.151}
\]
The second inequality includes ramified primes.
For the known six-dimensional \(X\), its archimedean
parameter in \(\operatorname{End}X\) has dimensions
\(20\) and \(16\) for the trivial and sign characters,
respectively.  At the specified prime five,
\[
 E_{\operatorname{End}X,5}(\sigma)
       -\frac3{10}E_{\operatorname{Reg},5}(\sigma)
       =\frac{4(\log5)5^{-\sigma}}{1-5^{-2\sigma}}.
 \tag{MA-RB.152}
\]

For the separate score assertion, let
\[
 \begin{gathered}
 (U_1,\ldots,U_5)=(W,W\eta,V,V\eta,X),\\
 a_i\ge0,\quad a\ne0,\qquad
 p=\sum_i a_i^2,\quad
 d=5a_1+5a_2+4a_3+4a_4+6a_5,\quad
 m=a^{\mathsf T}Ma ,
 \end{gathered}
 \tag{MA-RB.153}
\]
where \(M\) is exactly (MA-RB.65).  Define the positive-part
score by
\[
 {\cal S}(a)=
       \frac{\max\{\sqrt m-\sqrt p,\,0\}^{\,2}}{d^2}.
 \tag{MA-RB.154}
\]
Then every such choice satisfies the strict bound
\[
                         {\cal S}(a)<\frac4{120}.
 \tag{MA-RB.155}
\]
This is the score obtained from the guaranteed multiplicity
when the unknown comparison order is bounded only by \(k\ge1\).
For the honest Artin subfamily \((V,V\eta,X)\), the conductor
bound in (MA-RB.151) applies to every integer Artin direct
sum, whose unramified factors match the corresponding
automorphic isobaric sum.
For weights involving the automorphic comparison \(A\),
no corresponding Artin-conductor bound is asserted.

\begin{proof}
\emph{The finite inertia and conductor comparisons.}
Average a positive definite Hermitian form over \(S_5\),
so that \(U\) is unitary.  The character of its endomorphism
representation is therefore
\[
             \chi_{\operatorname{End}U}(g)=|\chi_U(g)|^2.
 \tag{MA-RB.156}
\]
For any subgroup \(I\subset S_5\), character averaging gives
\[
 \dim(\operatorname{End}U)^I
   =\frac1{|I|}\sum_{i\in I}|\chi_U(i)|^2
   \ge\frac{d^2}{|I|}.
\]
The regular representation has \(I\)-invariant dimension
\(120/|I|\).  Subtracting both invariant dimensions from
the respective total dimensions yields
\[
 \operatorname{codim}(\operatorname{End}U)^I
      \le\frac{d^2}{120}\operatorname{codim}
                                      (\operatorname{Reg})^I.
 \tag{MA-RB.157}
\]
The finite Artin-conductor exponent at \(p\), in lower
ramification numbering, is the positive weighted sum
\[
 a_p(T)=\sum_{r\ge0}
          \frac{|I_{p,r}|}{|I_{p,0}|}
                        \operatorname{codim}T^{I_{p,r}},
       \qquad I_{p,0}=I_p.
 \tag{MA-RB.158}
\]
Only finitely many summands are nonzero.  Applying
(MA-RB.157) to each actual ramification group proves
\(a_p(\operatorname{End}U)\le(d^2/120)a_p(\operatorname{Reg})\).
The regular conductor--discriminant identity is
\(N(\operatorname{Reg})=\Delta\).  Multiplication over
primes, or summation after logarithms, proves the first
line of (MA-RB.151).  This uses no classification of
inertia subgroups.

\emph{The actual ramified Frobenius trace.}
The averaging projector
\(|I_p|^{-1}\sum_{i\in I_p}T(i)\) commutes with
\(T(\phi_p^j)\), since \(\phi_p\) normalizes inertia.
Taking the trace of their product gives
\[
 t_{\operatorname{End}U,p}(j)
   =\frac1{|I_p|}\sum_{i\in I_p}
                        |\chi_U(i\phi_p^j)|^2 .
 \tag{MA-RB.159}
\]
If \(\phi_p^j\notin I_p\), the coset contains no
identity element, so the analogous regular trace is
zero and the displayed sum is nonnegative.
If \(\phi_p^j\in I_p\), the coset is \(I_p\) and
contains the identity.  The sum is then at least
\(d^2/|I_p|\), while the regular trace is \(120/|I_p|\).
This proves the second line of (MA-RB.151).
It also proves that every summand defining
\(E_{\operatorname{End}U}\) is nonnegative.
Absolute convergence for \(\sigma>1\) follows from
the fixed-dimensional finite-image bound, so the
termwise inequalities may be summed there.

For \(X\), the already checked character table gives
\(\chi_X(2B)=-2\).  Consequently
\(\chi_{\operatorname{End}X}(2B)=4\), and solving
\(d_++d_-=36\), \(d_+-d_-=4\) gives \(d_+=20,d_-=16\).
At five inertia is trivial and Frobenius has order two.
Its odd powers have endomorphism trace \(4\) and regular
trace \(0\); its even powers have the respective traces
\(36\) and \(120\).  Summing the odd powers proves
(MA-RB.152).

\emph{The exact character bound for the complete weight cone.}
The known class order is
\((1,2A,2B,3A,6A,4A,5A)\).  The genuine character
\(\Theta=V\eta+W\) has values
\[
                   (9,-1,1,0,2,-1,-1).
 \tag{MA-RB.160}
\]
These are read from the previously checked table, not
from a new representation certificate.  In particular
\(\Theta(g)\le2\) for every nonidentity element.
For real weights in (MA-RB.153), use the class function
\(\chi_U=\sum_i a_i\chi_{U_i}\).
Orthogonality of the five distinct irreducible
characters gives
\(\langle|\chi_U|^2,\mathbf1\rangle=p\), and
\(\chi_U(1)=d\).
The exact discrepancy matrix in (MA-RB.62)--(MA-RB.64)
gives
\[
 \begin{split}
 m&=\langle|\chi_U|^2,\Theta\rangle
                                -2(a_1-a_2)^2\\
  &\le\frac{9d^2}{120}
                    +2\left(p-\frac{d^2}{120}\right)
   =2p+\frac7{120}d^2 .
 \end{split}
 \tag{MA-RB.161}
\]
The inequality is valid because \(|\chi_U(g)|^2\ge0\)
on every class; no positivity of the character
\(\Theta\) itself is assumed.

Put \(x=p/d^2\).  Cauchy--Schwarz gives
\[
 d^2\le
  (5^2+5^2+4^2+4^2+6^2)\sum_i a_i^2=118p.
\]
Nonnegativity of the weights gives
\(d\ge4\sum_i a_i\) and
\((\sum_i a_i)^2\ge p\).  Thus
\[
                         \frac1{118}\le x\le\frac1{16}.
 \tag{MA-RB.162}
\]
If \(m\le p\), the score is zero and (MA-RB.155)
follows.  Otherwise (MA-RB.161) gives
\[
 {\cal S}(a)
   \le\bigl(\sqrt{2x+7/120}-\sqrt x\bigr)^2
                     <\frac1{30}.
 \tag{MA-RB.163}
\]
To verify the strict inequality without numerical
optimization, its positive square-root comparison is
\(\sqrt{2x+7/120}<\sqrt x+1/\sqrt{30}\).
After squaring and moving terms it is
\(x+1/40<2\sqrt{x/30}\).  Both sides are positive.
A further square makes it exactly
\[
                 (x-\tfrac1{120})(x-\tfrac3{40})<0.
\]
The whole interval in (MA-RB.162) lies strictly between
these roots.  This proves (MA-RB.155).
If only \(V,V\eta,X\) occur, Cauchy--Schwarz gives the
stronger lower bound \(x\ge1/68\), so the same strict
inequality applies.
\end{proof}

\paragraph{Why the positive part and the pole term matter.}
The elementary near-one model underlying the proposed
score is
\[
       \frac{m}{t+\delta}\le\frac p t+h,
          \qquad t>0,\quad \delta>0,\quad h>0,\quad p>0.
 \tag{MA-RB.164}
\]
Here \(p\) is the actual pole order, rather than a
dimension-scaled fraction of it.  Rearranging gives
\(\delta\ge mt/(p+ht)-t\).
The derivative of the right side is
\(mp/(p+ht)^2-1\).  If \(m\le p\), its supremum
over \(t>0\) is zero, approached as \(t\downarrow0\).
If \(m>p\), its maximum occurs at
\(t=(\sqrt{mp}-p)/h\) and equals
\[
                  \frac{(\sqrt m-\sqrt p)^2}{h}.
 \tag{MA-RB.165}
\]
This proves the positive-part convention in (MA-RB.154).
For integer isobaric multiplicities among the known
five distinct cuspidals, the self-Rankin pole order
is exactly \(\sum_i a_i^2=p\).  The established order
matrix guarantees multiplicity at least \(km\) at the
actual residual zero.  The certificate above concerns
the unconditional lower bound \(k\ge1\), not an
assertion that all actual orders equal this bound.

Under a degree-proportional conductor bound
\(\log N\le(d^2/120)\log\Delta\), the leading constant
in this model is \(240{\cal S}(a)/\log\Delta\).
The regular representation has the corresponding
multiplicity \(9\), pole order \(1\), and degree \(120\),
giving \(8/\log\Delta\).
For \(\operatorname{End}X\) alone they give
\[
 \frac{(\sqrt3-1)^2}{0.15\log\Delta}
                         <\frac8{\log\Delta}.
 \tag{MA-RB.166}
\]
This comparison is a statement about the displayed
model and its leading conductor score.  It is not
a substitution of that model for an exact zero sum.

\paragraph{The exact archimedean and Euler comparison.}
The additional terms can be kept visibly separate.
Let
\[
 \begin{gathered}
 P(\sigma)=\frac1\sigma+\frac1{\sigma-1},
 \quad
 g_0(\sigma)=\frac{\Gamma_{\mathbb R}'}{\Gamma_{\mathbb R}}
                                                   (\sigma),
 \quad
 g_1(\sigma)=\frac{\Gamma_{\mathbb R}'}{\Gamma_{\mathbb R}}
                                                   (\sigma+1),\\
 H_X^{\rm ex}(\sigma)
   =P(\sigma)+\tfrac12\log N(\operatorname{End}X)
                         +20g_0(\sigma)+16g_1(\sigma)
                         -E_{\operatorname{End}X}(\sigma),\\
 H_{\rm reg}^{\rm ex}(\sigma)
   =P(\sigma)+\tfrac12\log\Delta
                         +60g_0(\sigma)+60g_1(\sigma)
                         -E_{\operatorname{Reg}}(\sigma).
 \end{gathered}
 \tag{MA-RB.167}
\]
Here the \(X\) expression uses the actual Artin
function \(L(s,\operatorname{End}X)\), with all its
inertia-invariant factors.  Its indicated normalized
completion is
\(s(s-1)N(\operatorname{End}X)^{s/2}
\Gamma_{\mathbb R}(s)^{20}\Gamma_{\mathbb R}(s+1)^{16}
L(s,\operatorname{End}X)\).
For the regular representation replace the conductor
by \(\Delta\), both exponents by \(60\), and the finite
function by \(\zeta_L\).  Logarithmic differentiation
proves that (MA-RB.167) gives their logarithmic
derivatives on \(\sigma>1\).
The already checked partial Rankin identity agrees
with this Artin function away from the finite omitted
set.  No identification of ramified tensor-product
factors is inferred here from equality of scalar
standard \(L\)-factors alone.  Such an identification
is unnecessary for the finite lemma or the score.
Their exact difference is
\[
 \begin{split}
 H_X^{\rm ex}-\tfrac3{10}H_{\rm reg}^{\rm ex}
   ={}&\tfrac7{10}P
       +\tfrac12\bigl(\log N(\operatorname{End}X)
                                      -\tfrac3{10}\log\Delta\bigr)\\
     &+\psi(\sigma/2)-\psi((\sigma+1)/2)\\
     &-\bigl(E_{\operatorname{End}X}
                                  -\tfrac3{10}E_{\operatorname{Reg}}\bigr).
 \end{split}
 \tag{MA-RB.168}
\]
Indeed
\(\Gamma_{\mathbb R}(s)=\pi^{-s/2}\Gamma(s/2)\),
so \(2(g_0-g_1)=\psi(\sigma/2)-\psi((\sigma+1)/2)\).
The latter difference is strictly negative:
differentiating the gamma integral gives
\(\psi'(x)=\operatorname{Var}(\log T)>0\) for a
gamma-distributed random variable of shape \(x>0\).
The conductor and Euler differences in (MA-RB.168)
have the favorable signs established above.
The positive term \(7P/10\) is the unscaled pole cost.
No sign for the whole right side, or universal ordering
of the corresponding exact variational minima, follows
merely from these separate signs.

\paragraph{Source scope, motive meaning and remaining gate.}
Calegari's Theorem~1.1, in the supplied primary
\texttt{S5ArtinWeb.tex}, lines123--136, supplies the
modularity of \(V,V\eta,X\).  The already checked
Rankin theorem and MA-RB-006 supply the actual
isobaric pole and residual-order assertions.
The ramification-filtration expression (MA-RB.158)
is the classical definition of the Artin conductor;
the regular conductor--discriminant identity is the
inherited input recorded with Neukirch, VII
(11.7)--(11.9), printed pp.~533--534/PDF pp.~541--542,
in GRH\_NIGHT NT-C06 and its source ledger.
That locator is not claimed as a new reading of
the original monograph pages in this section.

In the finite-cover language already fixed by SGA~1,
Expos\'e~V, \S7, and Deligne's Artin-motive locators
recorded in MA-S5, \(\operatorname{End}U=U\otimes U^\vee\)
is an actual finite Artin object.  Its local factor
uses the actual inertia invariants and residue
Frobenius in (MA-RB.159).  The inequality against
a rational multiple of the regular character is
a scalar inequality; it is not a submotive assertion
or positivity of a space of analytic zeros.

The newer exceptional-zero statements checked for
this route are Harcos--Thorner (2025),
\emph{A new zero-free region for Rankin--Selberg
\(L\)-functions}, arXiv:2303.16889v6, Theorem~1.1,
printed p.~3, and Harcos--Thorner (2026),
\emph{Tatuzawa's theorem for
Rankin--Selberg \(L\)-functions},
arXiv:2508.10844v3, Theorem~1.1/Corollary~1.4,
p.~3, and Theorem~1.6, pp.~4--5.
Their exact hypotheses and conductor-dependent
regions were inspected.  No proof of coverage of
the residual interval, nor a fresh full audit of
those source proofs, is claimed.
In particular, a possible quadratic twist satisfies
the latter theorem's condition
\(\pi\otimes\chi^2=\pi\); odd degree does not
contradict that condition.

Thus the local lemma and the cone score are checked
scope results.  They do not exclude the actual
shared zero in (MA-RB.90), prove \((\mathcal R_B)\),
or construct a positive-excess field.
